% Generated draft: requires mathematical and source review.
\documentclass[11pt]{article}
\usepackage[margin=25mm]{geometry}
\usepackage{fontspec}
\setmainfont{Latin Modern Roman}
\usepackage{amsmath,amssymb,mathtools,unicode-math}
\setmathfont{Latin Modern Math}
\usepackage{enumitem,needspace,xurl,hyperref}
\setlength{\parindent}{0pt}
\setlength{\parskip}{5pt}
\setlength{\emergencystretch}{3em}
\newcommand{\N}{\mathbb{N}}
\newcommand{\Z}{\mathbb{Z}}
\newcommand{\Q}{\mathbb{Q}}
\newcommand{\R}{\mathbb{R}}
\newcommand{\C}{\mathbb{C}}
\newcommand{\F}{\mathbb{F}}
\newcommand{\A}{\mathbb{A}}
\newcommand{\PP}{\mathbb{P}}
\newcommand{\E}{\mathbb{E}}
\newcommand{\eps}{\varepsilon}
\newcommand{\dd}{\,\mathrm{d}}
\begin{document}
\section*{UnsolvedMath ID: 4900002}
\textbf{Source problem code:} catalogue entry
\subsection{Definitions and mathematical statement}
A classical Hamiltonian system with $f$ degrees of freedom is Liouville integrable when it has $f$ functionally independent integrals in involution on the regular region. The source considers quantum systems whose classical counterparts are integrable, but gives no quantization class or precise genericity/nondegeneracy assumptions.

For a quantum Hamiltonian with discrete ordered energies $E_1\leq E_2\leq\cdots\to\infty$, let $\overline N(E)$ be a specified smooth mean counting function, usually a Weyl approximation. Unfolding replaces the levels by $x_j=\overline N(E_j)$ so their mean density is one, and the nearest-neighbor gaps are $s_j=x_{j+1}-x_j$.

\textbf{Berry--Tabor conjecture.} For a generic quantum system with integrable classical counterpart, the unfolded high-energy level spacings have Poisson statistics. A precise nearest-neighbor consequence, under the above empirical averaging convention, would be
\[
 \frac1N\#\{1\leq j\leq N:s_j\leq s\}\longrightarrow1-e^{-s}
 \quad\text{for every }s\geq0.
\]
Full Poisson local-level statistics are stronger: after a specified high-energy local averaging or ensemble convention, counts in disjoint unfolded intervals should be independent Poisson random variables with means equal to interval lengths. The excerpt does not select a point-process convention or say whether only the gap distribution is intended. Poisson pair correlation alone does not state the full spacing assertion.

The intended space of systems and meaning of generic, degrees of freedom, allowed arithmetic resonances, multiplicities/symmetry sectors, unfolding, and energy-window averaging are not specified. Thus the conjecture is not a universal claim about every integrable example, nor is high energy silently replaced by a different semiclassical limit.

The supplied summary reports partial results for specific integrable models, mentioning two-point/pair correlation for almost all flat tori, inhomogeneous quadratic forms, and an Aharonov--Bohm ring under Diophantine assumptions. It leaves the generic full-spacing claim open. These restricted and weaker-statistic results have no supplied theorem citations and are not presented as proof of the general conjecture.
\subsection{Short English statement}
Generic quantum systems with integrable classical dynamics are predicted to have Poisson statistics after their high-energy levels are rescaled to unit mean density. In particular, nearest-neighbor gaps should have an exponential distribution. Genericity, spectral averaging, and the intended strength of Poisson statistics remain unspecified; reported pair-correlation results are weaker.
\subsection{Sources}
\begin{itemize}
\item Local UnsolvedMath registry, permanent ID 4900002, source code catalogue entry, ``Berry--Tabor conjecture''; proposers: Michael Berry and Michael Tabor. Supplied status partially solved, classification PARTIAL-PROGRESS, updated 2026-08-06.
\item Wikipedia, \emph{List of unsolved problems in mathematics}, Dynamical systems, current bullet 2 as extracted, \url{https://en.wikipedia.org/wiki/List_of_unsolved_problems_in_mathematics#Dynamical_systems}. Reported HTML extraction and access: 2026-07-29. The record says Wikipedia lists the item as unsolved but is not an authoritative resolution source. Rights: CC BY-SA 4.0; attribution and share-alike compatibility require release review.
\item The summary mentions flat tori, inhomogeneous quadratic forms, and an Aharonov--Bohm ring with Diophantine hypotheses, without publications or precise theorem statements. No original Berry--Tabor paper is supplied.
\end{itemize}
\end{document}
